\documentclass[11pt]{article}
\usepackage[utf8]{inputenc}
\usepackage{amsmath}
\usepackage{amsfonts}
\usepackage{amssymb}
\usepackage{bm}
\usepackage{geometry}
\usepackage{xcolor}
\usepackage{booktabs}

\geometry{margin=0.8in}

\title{Preprocessing Experimental Spike Recordings for PRISM-EM\\
\large Ver3c biological-data input preparation}
\author{Jan Balewski, NERSC}
\date{\today}

\begin{document}

\maketitle

\begin{abstract}
This document describes the current \texttt{prep\_bioexp3c.py}
preprocessing program for biological spike recordings. The script reads
raw experimental spike times and a per-channel metrics spreadsheet,
converts spike times into a binned spike-count matrix, filters channels by
firing-rate range, orders accepted neurons by firing rate, and writes two
companion files. The first file, \texttt{.spikes.npz}, contains only the
data needed by the EM fitter. The second file, \texttt{.bioExp.npz},
contains experimental metadata, MEA identifiers, metrics-table rows, and
other biological bookkeeping used by inspection, plotting, and downstream
interpretation.
\end{abstract}

\tableofcontents
\newpage

% ===============================================================
\section{Purpose and Scope}
\label{sec:scope}
% ===============================================================

\texttt{prep\_bioexp3c.py} converts one experimental recording session into
the same spike-file contract used elsewhere in the ver3c pipeline:
rows are time bins and columns are neurons. The preprocessing script does
not fit a model and does not estimate connectivity. Its job is to make the
experimental recording usable by the EM fitter while keeping
experiment-specific information outside the fitter input.

The key design choice is a two-file output:
\begin{itemize}
  \item \texttt{<dataName>.spikes.npz}: compact fitter input;
  \item \texttt{<dataName>.bioExp.npz}: biological/session metadata and
        channel annotations.
\end{itemize}
The files share the same stem and are written to the same directory.

% ===============================================================
\section{Command-Line Inputs}
\label{sec:cli}
% ===============================================================

The command-line arguments are:
\begin{itemize}
  \item \texttt{--expPath}: root directory containing the raw experimental
        session directory;
  \item \texttt{--sessionName}: session path below \texttt{--expPath};
  \item \texttt{--dataPath}: output directory for both generated
        \texttt{.npz} files;
  \item \texttt{--shortName}: optional output stem. If omitted, the script
        uses \texttt{<recording\_date>-<6hex>};
  \item \texttt{--inpFormat}: metrics-spreadsheet selector,
        \texttt{1} for \texttt{metrics\_curated.xlsx} and
        \texttt{2} for \texttt{quality\_metrics.xlsx};
  \item \texttt{--samp\_freq}: binning frequency in Hz, default
        \texttt{100};
  \item \texttt{--freqRange}: inclusive accepted firing-rate range in Hz,
        default \texttt{1 50};
  \item \texttt{--verb}: diagnostic verbosity.
\end{itemize}

The raw session is expected to contain
\[
\texttt{<expPath>/<sessionName>/spike\_times.npy}
\]
and one metrics spreadsheet:
\[
\texttt{metrics\_curated.xlsx}
\quad\hbox{or}\quad
\texttt{quality\_metrics.xlsx}.
\]

The session name is parsed as a slash-separated path whose fields include
cell line, recording date, chip ID, run number, and well number. These
values are stored in the biological metadata.

% ===============================================================
\section{Preprocessing Steps}
\label{sec:steps}
% ===============================================================

\subsection{Read Raw Spike Times}

The input \texttt{spike\_times.npy} is loaded as a Python dictionary keyed
by MEA feature/channel identifiers. The keys are sorted once, and this
sorted key order defines the initial raw channel order.

For each channel, spike times are multiplied by \texttt{--samp\_freq} and
cast to integer bin indices. Thus \texttt{--samp\_freq=100} corresponds to
\(\Delta t=0.01\,\mathrm{s}\). Repeated bin indices are retained and later
become spike counts larger than one in a bin.

The script computes:
\begin{itemize}
  \item the total number of time bins;
  \item the recording duration in seconds;
  \item one firing rate per raw channel.
\end{itemize}

\subsection{Filter by Firing Rate}

Channels are retained if their firing rate lies inside the inclusive range
\[
\texttt{freqRange[0]} \le r_i \le \texttt{freqRange[1]}.
\]
The metadata records how many channels were dropped below and above this
range. The remaining channels define the population used downstream.

\subsection{Sort Accepted Neurons by Firing Rate}

After rate filtering, accepted channels are sorted by firing rate in
ascending order. This frequency-sorted order becomes the primary neuron
index used in the saved spike matrix. The \texttt{.bioExp.npz} file stores
both the frequency-sort index and the inverse map so that downstream tools
can relate the saved columns back to the accepted-channel order.

\subsection{Build the Spike-Count Matrix}

For each accepted, frequency-sorted channel, the script forms a
length-\(T\) spike-count vector using \texttt{numpy.bincount}. These vectors
are stacked into
\[
Y \in \mathbb{N}_0^{T\times N},
\]
where rows are time bins and columns are accepted neurons.

The maximum pre-clipping count in any bin is recorded in the biological
metadata. The matrix written to the fitter-facing file is clipped to
\([0,255]\) and stored as \texttt{uint8}. This keeps the spike file compact
while preserving ordinary binned spike counts.

\subsection{Load and Align Metrics}

The metrics spreadsheet is loaded with \texttt{pandas}. For input format 1,
the file name is \texttt{metrics\_curated.xlsx}. For input format 2, the
file name is \texttt{quality\_metrics.xlsx}; location columns named
\texttt{location\_X}, \texttt{location\_Y}, and \texttt{location\_Z} are
renamed to \texttt{loc\_x}, \texttt{loc\_y}, and \texttt{loc\_z} when
present.

Rows are matched by \texttt{MEA\_idx} and reordered to match the final
frequency-sorted neuron order. Duplicate \texttt{MEA\_idx} rows are treated
as an error, and every retained channel must have a corresponding metrics
row.

\subsection{Compute Summary Statistics}

The script computes summary statistics from the accepted channels:
mean, standard deviation, median, minimum, and maximum firing rate, plus
the mean and standard deviation of the per-neuron Fano factor. These values
are stored in the biological metadata.

% ===============================================================
\section{Output File Pattern}
\label{sec:files}
% ===============================================================

Let \texttt{<dataName>} denote the output stem. It is equal to
\texttt{--shortName} when that argument is supplied; otherwise it is
generated as
\[
\texttt{<recording\_date>-<6hex>}.
\]

The script writes both output files directly into \texttt{--dataPath}:
\[
\texttt{<dataPath>/<dataName>.bioExp.npz},
\qquad
\texttt{<dataPath>/<dataName>.spikes.npz}.
\]

The order is deliberate: the \texttt{.spikes.npz} file is the direct input
to numerical fitters, while the \texttt{.bioExp.npz} file stores
experimental context that should not be required by the EM training code.

% ===============================================================
\section{Saved Contents}
\label{sec:saved}
% ===============================================================

\subsection{\texttt{<dataName>.spikes.npz}}

The fitter-facing data dictionary contains:
\begin{itemize}
  \item \texttt{spikes}: \texttt{uint8} array with shape \((T,N)\);
  \item \texttt{single\_rates}: floating-point array with shape \((N,)\),
        ordered exactly like the spike-matrix columns.
\end{itemize}

The metadata contains:
\begin{itemize}
  \item \texttt{time\_step\_sec}: bin width, equal to
        \(1/\texttt{samp\_freq}\);
  \item \texttt{provenance.experiment\_name}: the output stem
        \texttt{<dataName>};
  \item \texttt{poisson\_eta\_clip}: clipping value expected by the EM
        fitter, currently \texttt{5};
  \item \texttt{data\_type}: \texttt{bioExp};
  \item \texttt{num\_neurons}: number of accepted channels \(N\).
\end{itemize}

\subsection{\texttt{<dataName>.bioExp.npz}}

The biological data dictionary contains:
\begin{itemize}
  \item \texttt{neur\_freqIdx}: index vector that sorts accepted channels
        by firing rate;
  \item \texttt{neur\_revFreqIdx}: inverse map from accepted-channel order
        to frequency-sorted order;
  \item \texttt{single\_rates}: firing rates in the final saved neuron
        order;
  \item \texttt{MEA\_idx}: MEA identifiers reordered to match the final
        neuron order;
  \item \texttt{spike\_key}: raw \texttt{spike\_times.npy} keys reordered
        to match the final neuron order;
  \item \texttt{metrics\_curated}: metrics-spreadsheet rows reordered to
        match the final neuron order.
\end{itemize}

The biological metadata contains:
\begin{itemize}
  \item \texttt{bioexp}: raw-data provenance such as
        \texttt{raw\_bioexp\_path}, \texttt{session\_name},
        \texttt{inp\_format}, \texttt{cell\_line\_name},
        \texttt{recording\_date}, \texttt{chip\_ID}, \texttt{run\_num},
        \texttt{well\_num}, \texttt{sampling\_freq},
        \texttt{num\_feature}, \texttt{num\_time\_bin}, and
        \texttt{max\_time};
  \item \texttt{data\_selector}: selection parameters and results,
        including \texttt{freq\_range},
        \texttt{num\_drop\_neur\_lo\_hi\_freq}, \texttt{num\_chan}, and
        \texttt{max\_spike\_per\_bin};
  \item \texttt{hash}: random six-character hash used when
        \texttt{--shortName} is omitted;
  \item \texttt{short\_name}: output stem \texttt{<dataName>};
  \item \texttt{metrics\_curated\_columns}: spreadsheet column names after
        any input-format normalization;
  \item \texttt{rate\_summary}: firing-rate and Fano-factor summary
        statistics for the accepted population.
\end{itemize}

% ===============================================================
\section{Example Commands}
\label{sec:commands}
% ===============================================================

For inspection or plotting only, the output directory can be any directory
chosen by \texttt{--dataPath}. For use with the ver3c EM fitter, choose the
same \texttt{spikesData} layout used by \texttt{fitEM\_states\_ver3c.tex}:
\begin{verbatim}
basePath=/path/to/analysis

./prep_bioexp3c.py \
  --expPath /path/to/raw/experiments \
  --sessionName <cellLine>/<recordingDate>/<chipID>/<assay>/<run>/<well> \
  --dataPath $basePath/spikesData \
  --shortName <dataName> \
  --inpFormat 1 \
  --samp_freq 100 \
  --freqRange 1 50
\end{verbatim}

This produces:
\[
\texttt{\$basePath/spikesData/<dataName>.spikes.npz},
\qquad
\texttt{\$basePath/spikesData/<dataName>.bioExp.npz}.
\]

The \texttt{.spikes.npz} file is then the input selected by
\texttt{--basePath \$basePath --dataName <dataName>} in the EM training
workflow. See \texttt{fitEM\_states\_ver3c.tex} for the EM model, fitting
options, and fit-output file pattern.

\end{document}
